\section{Problem Statement}

Urban monitoring systems in modern smart cities are increasingly transitioning toward data-driven architectures; however, several critical gaps remain unresolved in high-frequency, real-time traffic anomaly detection pipelines. While current frameworks utilize \textbf{Kafka-based ingestion} to handle high-throughput traffic signals, there remains a persistent \textit{semantic gap between raw telemetry and actionable intelligence}.

In the context of the UrbanPulse system, traffic data is ingested from the VietMap API—capturing travel duration and congestion ratios—and processed through a \textbf{Medallion Lakehouse architecture} (comprising Bronze, Silver, and Gold layers). Despite this structured approach, interpreting anomalies in a streaming environment introduces several challenges:

\begin{itemize}
\item \textbf{Signal Inconsistency}: The system utilizes a \textbf{dual-signal detection framework} that combines statistical methods (Z-score) with machine learning models (Isolation Forest). Statistical detection often operates on real-time windows with restricted temporal context, whereas machine learning models rely on periodically retrained historical features, leading to potential misalignment in anomaly interpretation.
\item \textbf{Cold-start and Data Sparsity}: Effective Z-score computation requires a stable historical baseline of road segment activity. During the initial system bootstrap or in low-variance traffic conditions, the lack of robust baseline statistics results in null or unreliable anomaly signals.
\item \textbf{Model Staleness}: Machine learning models stored in the \textbf{MinIO-based model registry} are typically retrained on a batch schedule. This introduces a "staleness" factor, where the model may fail to adapt to sudden distribution shifts caused by unpredicted road incidents or temporary construction.
\item \textbf{High-Dimensional Latency Constraints}: Ensuring low-latency propagation from the \textbf{Redpanda broker} to the end-user dashboard via \textbf{Server-Sent Events (SSE)} requires a highly optimized serving layer that must reconcile high-frequency ingestion without introducing significant lag.
\end{itemize}

Therefore, the core problem addressed in this research is:
\begin{quote}
\textit{How to design a \textbf{streaming-first, low-latency traffic anomaly detection architecture} that ensures consistency between real-time signals and historical baselines while maintaining scalability through a unified lakehouse strategy?}
\end{quote}

\section{Scope}

This research focuses on the design and implementation of the "UrbanPulse" platform, an integrated system for real-time traffic monitoring and event detection within Ho Chi Minh City. The scope is defined as follows:

\begin{enumerate}
\item \textbf{Data Ingestion and Sources}
\begin{itemize}
\item Real-time traffic flow data is sourced exclusively from \textbf{VietMap APIs}, capturing travel duration, congestion ratios, and severity levels for a predefined set of key urban routes.
\item The data collection interval is configured at 5-minute increments to balance API constraints with real-time requirements.
\end{itemize}
\item \textbf{Technical and Architectural Framework}
\begin{itemize}
\item The system leverages a \textbf{Kafka messaging backbone} for high-throughput event propagation and fault-tolerant ingestion.
\item Data storage and lifecycle management are governed by a \textbf{Lakehouse pattern} utilizing \textbf{Apache Iceberg} for table formats, \textbf{Project Nessie} for catalog versioning, and \textbf{MinIO} as the underlying object store.
\item The processing pipeline is implemented in \textbf{Python}, utilizing micro-batching logic to transition data across the Medallion layers (Bronze $\rightarrow$ Silver $\rightarrow$ Gold).
\end{itemize}
\item \textbf{Analytical and Serving Boundary}
\begin{itemize}
\item Anomaly detection is limited to travel-time deviations and congestion surges using \textbf{Z-score statistics} and \textbf{Isolation Forest} models.
\item Results are served through a \textbf{FastAPI-based REST layer} and pushed to a real-time monitoring dashboard using \textbf{Server-Sent Events (SSE)}.
\end{itemize}
\end{enumerate}

\section{Objectives}

The primary goal of this thesis is to engineer a scalable, streaming-first platform for detecting traffic anomalies. Specific objectives include:

\begin{itemize}
\item \textbf{Architecture Realization}: Implement a robust Medallion-based data lakehouse that supports incremental processing of \textbf{Bronze} (raw), \textbf{Silver} (cleaned), and \textbf{Gold} (aggregated) traffic datasets.
\item \textbf{Low-Latency Propagation}: Achieve an end-to-end latency where anomaly signals are delivered to the serving layer within seconds of ingestion from the traffic API.
\item \textbf{Unified Detection Framework}: Develop a \textbf{dual-signal detector} that reconciles statistical Z-scores with Isolation Forest inference to minimize False Positive Rates (FPR) in traffic monitoring.
\item \textbf{Stream-to-Batch Consistency}: Ensure that historical traffic baselines stored in \textbf{Iceberg tables} are semantically consistent with the real-time features extracted during stream processing.
\item \textbf{Operational Observability}: Provide a monitoring interface that displays real-time congestion indicators, ingestion lag, and route-level anomaly insights.
\end{itemize}